% ============================================================
% §03a — GXL: GERT Expression Language (Normative)
% Author:  Edith — Spec Editor
% Created: 2026-06-05
% Source of truth: design/gert/grammar/gxl.ebnf (v1.0.0-draft)
% ============================================================
%
% RFC 2119 normative language used throughout:
%   MUST / MUST NOT  — absolute requirements for conforming implementations.
%   SHOULD / SHOULD NOT — recommended but not mandatory.
%   MAY              — optional; permitted but not required.
%
% Grammar production citations take the form: gxl.ebnf §<section>.<subsection>.
% Decision ledger citations take the form: decisions.md "<section-heading>" <key>.

% Local definition — remove if \TODO is promoted to the main preamble.
\providecommand{\TODO}[1]{\textbf{\textcolor{red}{[TODO:~#1]}}}

% ─────────────────────────────────────────────────────────────────────────────
\chapter{GXL\,---\,GERT Expression Language}
\label{sec:gxl}
% ─────────────────────────────────────────────────────────────────────────────

\begin{quote}
\noindent\textbf{Informative introduction.}
GXL (GERT Expression Language) is the GERT-native boolean expression language.
It is a small, deterministic, side-effect-free language designed to be
implementable from its grammar alone in any host language.
GXL is evaluated at every boolean-typed field in a runbook document:
\texttt{when}, \texttt{condition}, \texttt{until}, \texttt{iterate.until},
\texttt{branches[].condition}, \texttt{include.when}, and
\texttt{collector.fields[].when}.
Its design goals are portability across Go, C\#, and browser runtimes;
auditability by non-developer operators; and strict governance
(no implicit coercion, no side effects, no dynamic dispatch).
GXL replaces the prior \texttt{expr-lang/expr} evaluation surface;
see \S\ref{sec:gxl:migration} for the migration note.
\end{quote}

% ─────────────────────────────────────────────────────────────────────────────
\section{Normative Reference to the Grammar}
\label{sec:gxl:normref}
% ─────────────────────────────────────────────────────────────────────────────

The file \texttt{design/gert/grammar/gxl.ebnf} (version~1.0.0-draft) is the
authoritative grammar for GXL.  This chapter is normative prose \emph{about}
that grammar; it does not supersede it.  Where this prose and the grammar
conflict, \textbf{the grammar wins} and the prose is in error.  Conforming
implementations MUST derive their parsers from \texttt{gxl.ebnf}, not from
this chapter alone.

This grammar uses ISO EBNF with PEG alternation semantics:
ordered alternation (\texttt{|}) chooses the first matching alternative,
which eliminates all ambiguity
(\texttt{gxl.ebnf} \S1).

% ─────────────────────────────────────────────────────────────────────────────
\section{Lexical Structure}
\label{sec:gxl:lexical}
% ─────────────────────────────────────────────────────────────────────────────

The GXL lexer operates character-by-character over a Unicode string
(the field value as decoded from the YAML document).
Tokens are produced in the order they appear;
whitespace and comments are discarded before the parser sees any token.
All productions in this section correspond to \texttt{gxl.ebnf} \S2.

% ─────────────────────────────────────────────────────────────────────────────
\subsection{Whitespace}
\label{subsec:gxl:whitespace}

Whitespace characters --- U+0020~SPACE, U+0009~CHARACTER~TABULATION,
U+000A~LINE~FEED, and U+000D~CARRIAGE~RETURN --- are insignificant between
tokens.
Conforming lexers MUST silently discard all inter-token whitespace.
Whitespace inside string literals is \emph{not} discarded (see
\S\ref{subsec:gxl:strings}).
YAML block scalars (\texttt{|} and \texttt{>}) preserve newlines in the field
value; because newlines are GXL whitespace, multiline expressions are valid.

\begin{minted}{text}
WHITESPACE = { ' ' | '\t' | '\n' | '\r' } ;
\end{minted}

(Source: \texttt{gxl.ebnf} \S2.1.)

% ─────────────────────────────────────────────────────────────────────────────
\subsection{Comments}
\label{subsec:gxl:comments}

GXL supports shell-style single-line comments beginning with \texttt{\#}.
A comment extends from the \texttt{\#} character to (but not including)
the next newline.
Conforming lexers MUST discard comments in the same pass as whitespace.
The \texttt{\#} character is NOT supported inside string literals.

\begin{minted}{text}
LINE_COMMENT = '#' { ANY_CHAR_EXCEPT_NEWLINE } NEWLINE ;
\end{minted}

The \texttt{\#} comment character was chosen over \texttt{//}
for YAML-author familiarity.

(Source: \texttt{gxl.ebnf} \S2.1.)

% ─────────────────────────────────────────────────────────────────────────────
\subsection{Keywords}
\label{subsec:gxl:keywords}

The following identifiers are reserved as keywords and MUST NOT be used as
variable names.  Conforming parsers MUST reject any keyword used as an
identifier with error \texttt{GXL-PARSE-010}.

\begin{center}
\begin{tabular}{ll}
\hline
\textbf{Keyword} & \textbf{Role} \\
\hline
\texttt{true}   & Boolean literal \\
\texttt{false}  & Boolean literal \\
\texttt{null}   & Null literal \\
\texttt{and}    & Logical AND operator \\
\texttt{or}     & Logical OR operator \\
\texttt{not}    & Logical NOT operator \\
\texttt{len}    & Top-level length builtin \\
\texttt{now}    & Top-level UTC timestamp builtin \\
\texttt{str}    & Namespace prefix (string helpers) \\
\texttt{list}   & Namespace prefix (list helpers) \\
\texttt{regex}  & Namespace prefix (regex helpers) \\
\texttt{math}   & Reserved namespace prefix (v2; no methods in v1) \\
\hline
\end{tabular}
\end{center}

(Source: \texttt{gxl.ebnf} \S2.2.)

% ─────────────────────────────────────────────────────────────────────────────
\subsection{Identifiers}
\label{subsec:gxl:identifiers}

GXL identifiers are ASCII-only.  An identifier begins with a letter or
underscore (\texttt{[a-zA-Z\_]}) and continues with zero or more letters,
decimal digits, or underscores (\texttt{[a-zA-Z0-9\_]}).

\begin{minted}{text}
IDENT_START = 'a'..'z' | 'A'..'Z' | '_' ;
IDENT_CONT  = 'a'..'z' | 'A'..'Z' | '0'..'9' | '_' ;
RAW_IDENT   = IDENT_START { IDENT_CONT } ;
IDENT       = (* RAW_IDENT that is NOT a KEYWORD *) ;
\end{minted}

\noindent\textbf{Rationale for ASCII-only restriction.}
Variable names originate from capture path declarations and step IDs,
which are YAML keys and are ASCII by GERT authoring convention.
Unicode identifiers admit visual-spoofing attacks (Cyrillic look-alike
characters), which in runbook governance contexts constitute a compliance risk.

Conforming parsers MUST reject any keyword used in identifier position with
error \texttt{GXL-PARSE-010}.

(Source: \texttt{gxl.ebnf} \S2.3.)

% ─────────────────────────────────────────────────────────────────────────────
\subsection{String Literals}
\label{subsec:gxl:strings}

String literals are delimited by matching quotation marks.
Both double-quoted and single-quoted forms are supported.
The opening and closing delimiter MUST be the same character.
The alternate quote character requires no escaping inside a literal
(e.g., \texttt{"it's fine"} and \texttt{'say "hello"'} are both valid).

\begin{minted}{text}
STRING        = DQUOTE_STRING | SQUOTE_STRING ;
DQUOTE_STRING = '"' { DQUOTE_CHAR } '"' ;
SQUOTE_STRING = "'" { SQUOTE_CHAR } "'" ;
DQUOTE_CHAR   = ESCAPE_SEQ | (* any Unicode except '"', '\', LF, CR *) ;
SQUOTE_CHAR   = ESCAPE_SEQ | (* any Unicode except "'", '\', LF, CR *) ;
\end{minted}

\noindent Newlines (LF, CR) MUST NOT appear unescaped inside a string literal;
encountering end-of-line before the closing delimiter is error
\texttt{GXL-PARSE-002}.

The following escape sequences are recognized:

\begin{center}
\begin{tabular}{ll}
\hline
\textbf{Sequence} & \textbf{Meaning} \\
\hline
\texttt{\textbackslash\textbackslash} & U+005C REVERSE SOLIDUS \\
\texttt{\textbackslash"} & U+0022 QUOTATION MARK \\
\texttt{\textbackslash'} & U+0027 APOSTROPHE \\
\texttt{\textbackslash n} & U+000A LINE FEED \\
\texttt{\textbackslash r} & U+000D CARRIAGE RETURN \\
\texttt{\textbackslash t} & U+0009 CHARACTER TABULATION \\
\texttt{\textbackslash 0} & U+0000 NULL \\
\texttt{\textbackslash u}$XXXX$ & Unicode code point U+$XXXX$ (exactly 4 hex digits) \\
\texttt{\textbackslash U}$XXXXXXXX$ & Unicode code point U+$XXXXXXXX$ (exactly 8 hex digits) \\
\hline
\end{tabular}
\end{center}

Any other character following \texttt{\textbackslash} is error
\texttt{GXL-PARSE-004}.
Implementations MUST NOT invent additional escape sequences.

(Source: \texttt{gxl.ebnf} \S2.4.)

% ─────────────────────────────────────────────────────────────────────────────
\subsection{Number Literals}
\label{subsec:gxl:numbers}

\begin{minted}{text}
NUMBER  = INTEGER [ FRAC ] [ EXP ] ;
INTEGER = '0' | ( '1'..'9' { '0'..'9' } ) ;
FRAC    = '.' '0'..'9' { '0'..'9' } ;
EXP     = ( 'e' | 'E' ) [ '+' | '-' ] '0'..'9' { '0'..'9' } ;
\end{minted}

Conforming implementations MUST treat all GXL numbers as IEEE~754
double-precision float64.
Integer literals are representable exactly for values up to $2^{53}$;
values beyond this range are silently rounded to the nearest representable
double (this is not an error;
see open issue OI-GXL-02 in \texttt{gxl.ebnf} \S7).

The following number formats are supported:
integers (\texttt{42}),
decimal fractions (\texttt{3.14}),
and scientific notation (\texttt{1.5e10}, \texttt{2E-3}).

The following formats are \emph{not} supported and MUST be rejected at parse time:
hexadecimal (\texttt{0xFF}),
octal (\texttt{0o77}),
binary (\texttt{0b1010}),
leading-zero integers (\texttt{007} is error \texttt{GXL-PARSE-003}),
\texttt{NaN}, and
\texttt{Infinity}.
Negative number values are produced by the \texttt{UnaryMinus} rule
(\S\ref{subsec:gxl:unary}), not by the lexer;
the \texttt{NUMBER} token itself is always non-negative.

(Source: \texttt{gxl.ebnf} \S2.5.)

% ─────────────────────────────────────────────────────────────────────────────
\subsection{Operators and Punctuation}
\label{subsec:gxl:operators}

The following operator and punctuation tokens are recognized:

\begin{center}
\begin{tabular}{ll}
\hline
\textbf{Token(s)} & \textbf{Role} \\
\hline
\texttt{==}~\texttt{!=}~\texttt{<}~\texttt{<=}~\texttt{>}~\texttt{>=}
  & Comparison operators \\
\texttt{+}~\texttt{-} & Additive arithmetic operators \\
\texttt{*}~\texttt{/}~\texttt{\%} & Multiplicative arithmetic operators \\
\texttt{(}~\texttt{)} & Grouping delimiters \\
\texttt{[}~\texttt{]} & Array-index delimiters (in GDP paths) \\
\texttt{,}          & Argument separator in function calls \\
\texttt{.}          & Path segment separator (in GDP) and namespace separator \\
\hline
\end{tabular}
\end{center}

(Source: \texttt{gxl.ebnf} \S2.6.)

% ─────────────────────────────────────────────────────────────────────────────
\subsection{Forbidden Tokens}
\label{subsec:gxl:forbidden-tokens}

The following tokens MUST be rejected by conforming parsers with error
\texttt{GXL-PARSE-007}.
Their appearance in any expression field is unconditionally invalid,
regardless of context.

\begin{center}
\begin{tabular}{ll}
\hline
\textbf{Forbidden token or construct} & \textbf{Rationale} \\
\hline
\texttt{\&\&} \texttt{||} \texttt{!} & Go/C-style logical operators \\
\texttt{|} & Pipe operator \\
\texttt{?} \texttt{:} & Ternary operator tokens \\
\texttt{:=} \texttt{=} & Assignment operators \\
\texttt{in} (infix) & Membership test infix form \\
\texttt{contains} (infix) & Substring/membership infix form \\
\texttt{+} (on strings) & String concatenation (\texttt{+} is numeric only) \\
Array/object construction & \texttt{\{key:~value\}} and \texttt{[a,~b]} literal forms \\
Bare function calls & Any call of form \texttt{f()} not matching \texttt{len()} \\
Negative array index & \texttt{items[-1]} \\
Unary \texttt{+} & \texttt{+expr} \\
\hline
\end{tabular}
\end{center}

(Source: \texttt{gxl.ebnf} \S2.6, \S3 comments.)

% ─────────────────────────────────────────────────────────────────────────────
\section{Syntactic Structure}
\label{sec:gxl:syntactic}
% ─────────────────────────────────────────────────────────────────────────────

The GXL parser is a recursive-descent PEG parser.
There is no left-recursion in the grammar; the grammar is safe for PEG parsers
and LL parsers alike.

% ─────────────────────────────────────────────────────────────────────────────
\subsection{Operator Precedence and Associativity}
\label{subsec:gxl:precedence}

Operator precedence and associativity are encoded by the grammar's production
hierarchy.
Table~\ref{tab:gxl-precedence} lists all levels from lowest to highest binding.

\begin{table}[ht]
\centering
\begin{tabular}{clll}
\hline
\textbf{Level} & \textbf{Operator / Form} & \textbf{Associativity} & \textbf{Grammar rule} \\
\hline
1 (lowest) & \texttt{or}                        & Left  & \texttt{OrExpr}    \\
2          & \texttt{and}                       & Left  & \texttt{AndExpr}   \\
3          & \texttt{not} (unary prefix)        & Right & \texttt{NotExpr}   \\
4          & \texttt{==}~\texttt{!=}~\texttt{<}~\texttt{<=}~\texttt{>}~\texttt{>=}
                                                & \emph{Non-assoc.} & \texttt{Comparison} \\
5          & \texttt{+}~\texttt{-}              & Left  & \texttt{Addition}  \\
6          & \texttt{*}~\texttt{/}~\texttt{\%}  & Left  & \texttt{Multiply}  \\
7          & \texttt{-} (unary prefix)          & Right & \texttt{UnaryMinus} \\
8 (highest) & Literals, paths, calls, \texttt{()} & ---  & \texttt{Primary}   \\
\hline
\end{tabular}
\caption{GXL operator precedence and associativity (lowest to highest).
  Source: \texttt{gxl.ebnf} \S3 precedence table.}
\label{tab:gxl-precedence}
\end{table}

\noindent\textbf{Non-associativity of comparisons.}
Comparison operators are non-associative: chaining them is a parse error.
The expression \texttt{a < b < c} is error \texttt{GXL-PARSE-009};
authors MUST write \texttt{a < b and b < c} instead.

(Source: \texttt{gxl.ebnf} \S3.)

% ─────────────────────────────────────────────────────────────────────────────
\subsection{Complete Expression Grammar}
\label{subsec:gxl:grammar}

\begin{minted}{text}
Expression   = OrExpr ;
OrExpr       = AndExpr { "or" AndExpr } ;
AndExpr      = NotExpr { "and" NotExpr } ;
NotExpr      = "not" NotExpr | Comparison ;
Comparison   = Addition [ CompOp Addition ] ;
CompOp       = "==" | "!=" | "<=" | ">=" | "<" | ">" ;
Addition     = Multiply { ("+" | "-") Multiply } ;
Multiply     = UnaryMinus { ("*" | "/" | "%") UnaryMinus } ;
UnaryMinus   = "-" UnaryMinus | Primary ;
Primary      = LenCall
             | NamespaceCall
             | GDP
             | Literal
             | "(" Expression ")"
             ;
\end{minted}

(Source: \texttt{gxl.ebnf} \S3.)

% ─────────────────────────────────────────────────────────────────────────────
\subsection{Literal Expressions}
\label{subsec:gxl:literals}

\begin{minted}{text}
Literal      = STRING | NUMBER | BoolLiteral | NullLiteral ;
BoolLiteral  = "true" | "false" ;
NullLiteral  = "null" ;
\end{minted}

The four GXL literal types are:

\begin{itemize}
  \item \textbf{String}: single- or double-quoted Unicode text
    (see \S\ref{subsec:gxl:strings}).
  \item \textbf{Number}: IEEE~754 float64 value
    (see \S\ref{subsec:gxl:numbers}).
  \item \textbf{Boolean}: the keywords \texttt{true} and \texttt{false}.
  \item \textbf{Null}: the keyword \texttt{null},
    representing an absent or unset value.
\end{itemize}

(Source: \texttt{gxl.ebnf} \S3.)

% ─────────────────────────────────────────────────────────────────────────────
\subsection{GERT Dotted Path (GDP) Variable Access}
\label{subsec:gxl:gdp}

GDP (GERT Dotted Path) is the syntax for reading values from the run
variable context map.

\begin{minted}{text}
GDP = IDENT { "." IDENT | "[" INTEGER "]" } ;
\end{minted}

The first \texttt{IDENT} is the variable name; it MUST exist in the run
variable context at the time of evaluation.
Subsequent segments navigate into the resolved value:
a dotted segment selects a named field from an object,
and a bracketed integer selects an element from an array (zero-based).

\noindent Examples:

\begin{minted}{text}
hostname              -- top-level variable
config.region         -- field "region" of captured object "config"
items[0].metadata     -- index 0 of array "items", then field "metadata"
\end{minted}

For the normative error behavior of missing segments, out-of-bounds indices,
and non-array indexing, see \S\ref{sec:gxl:semantics:paths}.

(Source: \texttt{gxl.ebnf} \S3.)

% ─────────────────────────────────────────────────────────────────────────────
\subsection{Function Calls}
\label{subsec:gxl:calls}

GXL provides exactly three call forms:
the top-level \texttt{now} builtin, the top-level \texttt{len} builtin,
and namespace-qualified method calls.
All other function-call forms MUST be rejected as \texttt{GXL-PARSE-007}.

\subsubsection{now (top-level builtin)}
\label{subsubsec:gxl:now-call}

\begin{minted}{text}
NowCall = "now" "(" [ ArgList ] ")" ;
\end{minted}

\texttt{now} takes exactly zero arguments and returns a \texttt{string}
containing the current UTC timestamp in ISO-8601 format:
\texttt{YYYY-MM-DDTHH:MM:SSZ} (no fractional seconds; \texttt{Z} suffix
mandatory).  Any argument provided raises \texttt{GXL-TYPE-004} at parse
time (arity is statically known to be zero).

\texttt{now} is NOT a namespace; it MUST be written as \texttt{now()},
not \texttt{now.\ldots(\ldots)}.  A bare \texttt{now} without parentheses
cannot match \texttt{NowCall} (parentheses required) and cannot match a GDP
variable path (\texttt{now} is a reserved keyword, not \texttt{IDENT}); the
correct error is \texttt{GXL-PARSE-001} (unexpected token).

\paragraph{Return format.}
The returned string satisfies the grammar:

\begin{minted}{text}
year  "-" month "-" day "T" hour ":" minute ":" second "Z"

where year   = 4DIGIT
      month  = 2DIGIT  (* 01..12 *)
      day    = 2DIGIT  (* 01..31 *)
      hour   = 2DIGIT  (* 00..23 *)
      minute = 2DIGIT  (* 00..59 *)
      second = 2DIGIT  (* 00..60, allowing leap second *)
\end{minted}

Equivalently, the value MUST match the regular expression:
\texttt{\textasciicircum\textbackslash d\{4\}-\textbackslash d\{2\}-\textbackslash d\{2\}T\textbackslash d\{2\}:\textbackslash d\{2\}:\textbackslash d\{2\}Z\$}.

\paragraph{Determinism semantics.}
\texttt{now()} is \emph{per-evaluation}: each call site yields an
independent fresh timestamp at the moment of evaluation.  This matches
the semantics of Go's \texttt{time.Now()} and avoids the complexity of a
per-run snapshot registry.  Implementations that require a stable timestamp
across an entire run MUST capture the value once via a capture step and
reference the variable thereafter.

\paragraph{Portability reference implementations.}

\begin{itemize}
  \item Go: \texttt{time.Now().UTC().Format("2006-01-02T15:04:05Z")}
    (use the explicit layout string, not \texttt{time.RFC3339}, which
    may emit sub-second precision depending on the \texttt{time.Time}
    value).
  \item C\#: \texttt{DateTime.UtcNow.ToString("yyyy-MM-dd\textbackslash{}THH:mm:ss\textbackslash{}Z")}
  \item JS: \texttt{new Date().toISOString().replace(/\textbackslash.\textbackslash d+/, '')}
    (strips fractional seconds from the ISO string).
\end{itemize}

\paragraph{Use-case examples.}
\texttt{now()} is intended for evidence timestamps and audit fields where
the current wall-clock time must be embedded at evaluation time:

\begin{minted}{yaml}
# SOC2 audit completion timestamp
completion_date: "${now()}"

# Evidence record audit field
created_at: "${now()}"
\end{minted}

(Source: \texttt{gxl.ebnf} \S4.1b.)

\subsubsection{len (top-level builtin)}
\label{subsubsec:gxl:len-call}

\begin{minted}{text}
LenCall = "len" "(" Expression ")" ;
\end{minted}

\texttt{len} accepts exactly one argument of type \texttt{string} or
\texttt{list} and returns a \texttt{number}.
A call with zero or more than one argument is error \texttt{GXL-TYPE-004}.
\texttt{len} is NOT a namespace; it MUST be written as \texttt{len(\ldots)},
not \texttt{len.\ldots(\ldots)}.

\subsubsection{Namespace method calls}
\label{subsubsec:gxl:ns-call}

\begin{minted}{text}
NamespaceCall = Namespace "." Method "(" ArgList ")" ;
Namespace     = "str" | "list" | "regex" ;
Method        = IDENT ;
ArgList       = [ Expression { "," Expression } ] ;
\end{minted}

The namespace prefix MUST be one of \texttt{str}, \texttt{list}, or
\texttt{regex}.
Referencing an unknown namespace is error \texttt{GXL-PARSE-005}.
Referencing an unknown method within a valid namespace is error
\texttt{GXL-PARSE-006} (raised at parse time; the stdlib function set is
closed per \S1.6 and method names are statically validated).
Referencing the \texttt{math} namespace in v1 is error \texttt{GXL-PARSE-005}
(\texttt{math} is reserved for v2; no methods are defined in this version).

\paragraph{Note on \texttt{str.foo} without parentheses.}
The expression \texttt{str.foo} (no call parentheses) does not match
\texttt{NamespaceCall} (which requires \texttt{(ArgList)}) and cannot match a
GDP variable path (\texttt{str} is a keyword, not \texttt{IDENT}).
Under PEG ordered alternation all \texttt{Primary} alternatives fail;
the correct error is \texttt{GXL-PARSE-001} (unexpected token), not
\texttt{GXL-PARSE-006} or \texttt{GXL-PARSE-010}.
See conformance vector \texttt{TV-GXL-PARSE-083}.

(Source: \texttt{gxl.ebnf} \S3.)

% ─────────────────────────────────────────────────────────────────────────────
\subsection{Unary Operators}
\label{subsec:gxl:unary}

GXL defines two unary operators:

\begin{itemize}
  \item \textbf{Logical NOT} (\texttt{not}):
    Precedence level~3 (see Table~\ref{tab:gxl-precedence}).
    Requires a \texttt{bool} operand; applying it to any other type is
    error \texttt{GXL-TYPE-002}.
    \texttt{not} binds tighter than \texttt{and}/\texttt{or}:
    \texttt{not a and b} parses as \texttt{(not a) and b}.

  \item \textbf{Unary minus} (\texttt{-}):
    Precedence level~7.
    Requires a \texttt{number} operand; applying it to any other type is
    error \texttt{GXL-TYPE-003}.
    Unary \texttt{+} is not supported and MUST be rejected as
    \texttt{GXL-PARSE-007}.
\end{itemize}

(Source: \texttt{gxl.ebnf} \S3.)

% ─────────────────────────────────────────────────────────────────────────────
\subsection{Binary Operators}
\label{subsec:gxl:binary}

\begin{itemize}
  \item \textbf{Logical OR / AND}:
    Left-associative (levels~1 and~2 respectively).
    Both operators short-circuit; see \S\ref{sec:gxl:semantics:short-circuit}.

  \item \textbf{Comparison operators}
    (\texttt{==}, \texttt{!=}, \texttt{<}, \texttt{<=}, \texttt{>}, \texttt{>=}):
    Non-associative (level~4).
    Strict typing rules apply; see \S\ref{sec:gxl:semantics:types}.

  \item \textbf{Arithmetic operators}
    (\texttt{+}, \texttt{-}, \texttt{*}, \texttt{/}, \texttt{\%}):
    Addition and subtraction at level~5 (left-associative);
    multiplication, division, and modulo at level~6 (left-associative).
    The \texttt{+} operator is numeric addition only;
    string concatenation is NOT supported and MUST be rejected as
    \texttt{GXL-PARSE-007}.
    Division-by-zero and modulo-by-zero rules are in
    \S\ref{sec:gxl:semantics:arithmetic}.
\end{itemize}

(Source: \texttt{gxl.ebnf} \S3.)

% ─────────────────────────────────────────────────────────────────────────────
\section{Standard Library Namespace Functions}
\label{sec:gxl:stdlib}
% ─────────────────────────────────────────────────────────────────────────────

All function calls in GXL MUST go through one of the top-level builtins
(\texttt{len}, \texttt{now}) or one of the permitted namespaces
(\texttt{str}, \texttt{list}, \texttt{regex}).
Method names use camelCase
(decisions.md ``GXL/GIS/GCP --- Open question resolutions''~Q3).
The type notation used in the signatures below is:
\texttt{string}, \texttt{number}, \texttt{bool}, \texttt{null},
\texttt{list}, and \texttt{scalar} (meaning \texttt{string}, \texttt{number},
or \texttt{bool}).

(Source: \texttt{gxl.ebnf} \S4.)

% ─────────────────────────────────────────────────────────────────────────────
\subsection{\texttt{len} --- Top-Level Builtin}
\label{subsec:gxl:stdlib:len}

\begin{center}
\begin{tabular}{p{5.5cm}p{1.5cm}p{6cm}}
\hline
\textbf{Signature} & \textbf{Return} & \textbf{Description} \\
\hline
\texttt{len(s:\ string)} & \texttt{number}
  & UTF-8 character count (code points, not bytes).
    Empty string returns~0. \\
\texttt{len(l:\ list)} & \texttt{number}
  & Element count. Empty list returns~0. \\
\hline
\end{tabular}
\end{center}

A \texttt{null} argument to \texttt{len} MUST raise \texttt{GXL-TYPE-003}.

(Source: \texttt{gxl.ebnf} \S4.1.)

% ─────────────────────────────────────────────────────────────────────────────
\subsection{\texttt{now} --- Top-Level Builtin}
\label{subsec:gxl:stdlib:now}

\begin{center}
\begin{tabular}{p{5.5cm}p{1.5cm}p{6cm}}
\hline
\textbf{Signature} & \textbf{Return} & \textbf{Description} \\
\hline
\texttt{now()} & \texttt{string}
  & Current UTC timestamp in ISO-8601 format
    \texttt{YYYY-MM-DDTHH:MM:SSZ}.
    No fractional seconds; \texttt{Z} suffix always present. \\
\hline
\end{tabular}
\end{center}

\texttt{now()} takes exactly zero arguments.
Passing any argument raises \texttt{GXL-TYPE-004} at parse time.

\paragraph{Return format.}
The returned value MUST match the regular expression
\texttt{\textasciicircum\textbackslash d\{4\}-\textbackslash d\{2\}-\textbackslash d\{2\}T\textbackslash d\{2\}:\textbackslash d\{2\}:\textbackslash d\{2\}Z\$}.
The format is a strict subset of RFC~3339 / ISO-8601.

\paragraph{Determinism.}
\texttt{now()} is \emph{per-evaluation}: each occurrence of \texttt{now()}
in an expression is evaluated independently and may return a different
timestamp if sufficient time elapses between evaluations.
Implementations that require a stable snapshot value across a single run
MUST capture the result once via a capture step and reference the captured
variable thereafter.
Rationale: per-evaluation semantics match the natural behaviour of
\texttt{time.Now()} in Go and require no per-run state registry.

\paragraph{Portability.}
Conforming implementations across execution environments MUST produce
output in the exact format specified above:

\begin{itemize}
  \item Go: \texttt{time.Now().UTC().Format("2006-01-02T15:04:05Z")}
  \item C\#: \texttt{DateTime.UtcNow.ToString("yyyy-MM-dd\textbackslash{}THH:mm:ss\textbackslash{}Z")}
  \item JS: \texttt{new Date().toISOString().replace(/\textbackslash.\textbackslash d+/, '')}
\end{itemize}

\paragraph{Use cases.}
Evidence timestamps, audit-trail fields, and SOC2 completion records
where the current wall-clock time must be embedded at evaluation time:
\texttt{completion\_date: "\$\{now()\}"}, \texttt{created\_at: "\$\{now()\}"}.

(Source: \texttt{gxl.ebnf} \S4.1b.)

% ─────────────────────────────────────────────────────────────────────────────
\subsection{\texttt{str} Namespace}
\label{subsec:gxl:stdlib:str}

\begin{center}
\begin{tabular}{p{5.5cm}p{1.5cm}p{6cm}}
\hline
\textbf{Signature} & \textbf{Return} & \textbf{Description} \\
\hline
\texttt{str.contains(s,\ substr:\ string)} & \texttt{bool}
  & True if \texttt{s} contains \texttt{substr} as a substring.
    Empty \texttt{substr} always returns \texttt{true}. \\
\texttt{str.startsWith(s,\ prefix:\ string)} & \texttt{bool}
  & True if \texttt{s} begins with \texttt{prefix}.
    Empty \texttt{prefix} always returns \texttt{true}. \\
\texttt{str.endsWith(s,\ suffix:\ string)} & \texttt{bool}
  & True if \texttt{s} ends with \texttt{suffix}.
    Empty \texttt{suffix} always returns \texttt{true}. \\
\texttt{str.toLower(s:\ string)} & \texttt{string}
  & Unicode full case-fold to lowercase. \\
\texttt{str.toUpper(s:\ string)} & \texttt{string}
  & Unicode full case-fold to uppercase. \\
\texttt{str.trim(s:\ string)} & \texttt{string}
  & Remove leading and trailing Unicode whitespace
    (U+0009, U+000A, U+000D, U+0020, and Unicode Zs category). \\
\texttt{str.trimPrefix(s,\ prefix:\ string)} & \texttt{string}
  & Remove \texttt{prefix} from \texttt{s} if present; otherwise return \texttt{s}. \\
\texttt{str.trimSuffix(s,\ suffix:\ string)} & \texttt{string}
  & Remove \texttt{suffix} from \texttt{s} if present; otherwise return \texttt{s}. \\
\texttt{str.length(s:\ string)} & \texttt{number}
  & Alias for \texttt{len(s)}. Provided for namespace consistency. \\
\hline
\end{tabular}
\end{center}

A \texttt{null} argument to any \texttt{str.*} function MUST raise
\texttt{GXL-TYPE-003}.

(Source: \texttt{gxl.ebnf} \S4.2.)

% ─────────────────────────────────────────────────────────────────────────────
\subsection{\texttt{list} Namespace}
\label{subsec:gxl:stdlib:list}

\begin{center}
\begin{tabular}{p{5.5cm}p{1.5cm}p{6cm}}
\hline
\textbf{Signature} & \textbf{Return} & \textbf{Description} \\
\hline
\texttt{list.contains(l:\ list,\ item:\ scalar)} & \texttt{bool}
  & True if \texttt{l} contains an element strictly equal to \texttt{item}
    (same type, same value).
    A cross-type \texttt{item} raises \texttt{GXL-TYPE-003}.
    \texttt{null}~\texttt{item} returns \texttt{false}
    (null is never in a list). \\
\texttt{list.indexOf(l:\ list,\ item:\ scalar)} & \texttt{number}
  & Zero-based index of the first element equal to \texttt{item},
    or $-1$ if not found.
    Cross-type \texttt{item} raises \texttt{GXL-TYPE-003}.
    \texttt{null}~\texttt{item} always returns $-1$. \\
\texttt{list.length(l:\ list)} & \texttt{number}
  & Alias for \texttt{len(l)}. Provided for namespace consistency. \\
\hline
\end{tabular}
\end{center}

A \texttt{null} list argument (\texttt{l}) to any \texttt{list.*} function
MUST raise \texttt{GXL-TYPE-003}.

(Source: \texttt{gxl.ebnf} \S4.3.)

% ─────────────────────────────────────────────────────────────────────────────
\subsection{\texttt{regex} Namespace}
\label{subsec:gxl:stdlib:regex}

\begin{center}
\begin{tabular}{p{5.5cm}p{1.5cm}p{6cm}}
\hline
\textbf{Signature} & \textbf{Return} & \textbf{Description} \\
\hline
\texttt{regex.match(s,\ pattern:\ string)} & \texttt{bool}
  & True if \texttt{s} matches the RE2 regular expression
    \texttt{pattern}.
    Pattern MUST be valid RE2 syntax; no PCRE, no backreferences.
    An invalid pattern raises \texttt{GXL-EVAL-003} at evaluation time.
    Either argument \texttt{null} raises \texttt{GXL-TYPE-003}. \\
\hline
\end{tabular}
\end{center}

(Source: \texttt{gxl.ebnf} \S4.4.)

% ─────────────────────────────────────────────────────────────────────────────
\subsection{\texttt{math} Namespace (Reserved --- v2)}
\label{subsec:gxl:stdlib:math}

The \texttt{math} namespace is reserved for a future version (v2).
No methods are defined in this version.
Any expression referencing \texttt{math.*} in a v1 context MUST be rejected
with error \texttt{GXL-PARSE-005}.

(Source: \texttt{gxl.ebnf} \S4.5.)

% ─────────────────────────────────────────────────────────────────────────────
\section{Semantics}
\label{sec:gxl:semantics}
% ─────────────────────────────────────────────────────────────────────────────

GXL expressions are \emph{pure}: they produce no I/O, perform no variable
mutation, and invoke no dynamic dispatch.
The same expression evaluated twice against the same context map MUST produce
the same result.

(Source: \texttt{gxl.ebnf} \S5.6.)

% ─────────────────────────────────────────────────────────────────────────────
\subsection{Type System}
\label{sec:gxl:semantics:types}
% ─────────────────────────────────────────────────────────────────────────────

GXL uses the GERT Portable JSON Value Model (PJVM), which defines six value
types:

\begin{center}
\begin{tabular}{lp{10cm}}
\hline
\textbf{Type}   & \textbf{Description} \\
\hline
\texttt{null}   & Absent or unset value.
                  Represented by the literal \texttt{null}. \\
\texttt{bool}   & Boolean value: \texttt{true} or \texttt{false}. \\
\texttt{number} & IEEE~754 double-precision float64. \\
\texttt{string} & Immutable sequence of Unicode code points. \\
\texttt{list}   & Ordered sequence of PJVM values.
                  Lists are accessible only through \texttt{list.*}
                  and \texttt{len()}; direct indexing via \texttt{[N]}
                  in GDP paths is supported for captured subtrees. \\
\texttt{object} & Key-value map of string keys to PJVM values.
                  Accessible via GDP field segments (\texttt{.field}). \\
\hline
\end{tabular}
\end{center}

\noindent\textbf{No implicit coercion.}
GXL performs no implicit type conversions.
Comparing values of incompatible types MUST raise \texttt{GXL-TYPE-001}
at evaluation time (or at plan time where statically inferrable).
Authors MUST NOT rely on implicit truthiness:
a bare string or number identifier used where a \texttt{bool} is required
MUST raise \texttt{GXL-TYPE-002}.

\noindent The following null-comparison rules apply
(\texttt{gxl.ebnf} \S5.2):

\begin{itemize}
  \item \texttt{null == null} evaluates to \texttt{true} (no error).
  \item \texttt{null != null} evaluates to \texttt{false} (no error).
  \item \texttt{x == null} evaluates to \texttt{true} if \texttt{x} is null,
    \texttt{false} otherwise (no error, regardless of \texttt{x}'s type).
  \item \texttt{x != null} evaluates to \texttt{true} if \texttt{x} is not null,
    \texttt{false} otherwise (no error).
  \item Using \texttt{null} with any ordered comparison operator
    (\texttt{<}, \texttt{<=}, \texttt{>}, \texttt{>=}) MUST raise
    \texttt{GXL-EVAL-004}.
  \item \texttt{not null} MUST raise \texttt{GXL-TYPE-002}.
  \item \texttt{null} in arithmetic position MUST raise \texttt{GXL-TYPE-003}.
\end{itemize}

(Source: \texttt{gxl.ebnf} \S5.2, \S5.5.)

% ─────────────────────────────────────────────────────────────────────────────
\subsection{Short-Circuit Evaluation}
\label{sec:gxl:semantics:short-circuit}
% ─────────────────────────────────────────────────────────────────────────────

\noindent\textbf{Source:} \texttt{gxl.ebnf} \S5.1;
decisions.md ``GXL/GIS/GCP --- Open question resolutions''~OQ1.

GXL \texttt{and} and \texttt{or} operators MUST use lazy (short-circuit)
evaluation.

\begin{itemize}
  \item \textbf{OrExpr}: The evaluator MUST evaluate the left operand first.
    If the left operand is \texttt{true}, the evaluator MUST return
    \texttt{true} immediately without evaluating the right operand.
    If the left operand is \texttt{false}, the evaluator MUST then evaluate
    and return the right operand.

  \item \textbf{AndExpr}: The evaluator MUST evaluate the left operand first.
    If the left operand is \texttt{false}, the evaluator MUST return
    \texttt{false} immediately without evaluating the right operand.
    If the left operand is \texttt{true}, the evaluator MUST then evaluate
    and return the right operand.
\end{itemize}

As a direct consequence of short-circuit evaluation: a missing variable or
type error in the right-hand operand of \texttt{and}/\texttt{or} MUST NOT
raise an error if the left-hand operand short-circuits.
This is intentional; it enables defensive guards such as:

\begin{minted}{text}
has_result and result.status == "ok"
\end{minted}

\noindent where \texttt{has\_result} prevents evaluation of
\texttt{result.status} when the capture did not run.

When multiple operands are chained, short-circuit applies at each pairwise
step from left to right:

\begin{minted}{text}
a and b and c   -- evaluates as: (a and b) and c
a or  b or  c   -- evaluates as: (a or  b) or  c
\end{minted}

Conformance test category: \texttt{TV-GXL-EVAL} (short-circuit vectors).

% ─────────────────────────────────────────────────────────────────────────────
\subsection{Path Access (GDP)}
\label{sec:gxl:semantics:paths}
% ─────────────────────────────────────────────────────────────────────────────

A GDP expression (\S\ref{subsec:gxl:gdp}) is evaluated against the run
variable context map as follows:

\begin{enumerate}
  \item The first \texttt{IDENT} is looked up in the context map.
    If not found, the evaluator MUST raise \texttt{GXL-PATH-001}.
  \item Each subsequent dotted segment (\texttt{.field}) is applied to the
    current value.
    If the current value is not an object, or if the field name is not present
    in the object, the evaluator MUST raise \texttt{GXL-PATH-001}.
  \item Each subsequent indexed segment (\texttt{[N]}) is applied to the
    current value.
    If the current value is not a list, the evaluator MUST raise
    \texttt{GXL-PATH-003}.
    If the index \texttt{N} is greater than or equal to the list length,
    the evaluator MUST raise \texttt{GXL-PATH-002}.
\end{enumerate}

Silent substitution of empty strings or null values for missing paths is
\textbf{forbidden}.
Every missing reference is a hard error.
Authors MUST use \texttt{when:} guards (exploiting short-circuit evaluation,
\S\ref{sec:gxl:semantics:short-circuit}) or capture defaults
(GCP grammar, \S\ref{sec:gxl:semantics:paths}) to handle optional values.

A bare GDP path in \emph{boolean position} (as the top-level expression,
or directly following \texttt{and}/\texttt{or}/\texttt{not}) MUST resolve
to a value of type \texttt{bool}.
A non-\texttt{bool} value in boolean position MUST raise
\texttt{GXL-TYPE-002}.

(Source: \texttt{gxl.ebnf} \S3, \S5.3.)

% ─────────────────────────────────────────────────────────────────────────────
\subsection{Arithmetic}
\label{sec:gxl:semantics:arithmetic}
% ─────────────────────────────────────────────────────────────────────────────

All arithmetic operations (\texttt{+}, \texttt{-}, \texttt{*},
\texttt{/}, \texttt{\%}) require both operands to be of type \texttt{number}.
Applying arithmetic to a non-\texttt{number} operand MUST raise
\texttt{GXL-TYPE-003}.
All arithmetic is performed in IEEE~754 double-precision float64.
Integer values are representable exactly up to $2^{53}$;
larger values incur silent rounding (not an error).

\noindent\textbf{Division and modulo by zero.}
The expressions \texttt{x / 0} and \texttt{x \% 0} MUST raise
\texttt{GXL-EVAL-002} at evaluation time.
No IEEE~754 Infinity or NaN value is ever produced as a GXL result.
Integer-like division uses float64 semantics: \texttt{5 / 2 = 2.5}.

\begin{quote}
\noindent\textit{Non-normative note.}
The \texttt{\%} operator follows IEEE~754 floating-point remainder
(\texttt{fmod}) semantics: the result has the sign of the dividend.
For example, \texttt{-7 \% 3 = -1.0}.
This matches Go's behavior but differs from Python's modulo,
which always returns a non-negative result when the divisor is positive.
Conformance corpus MUST include negative-operand modulo test vectors
before this behavior can be considered fully specified.
(Open issue OI-GXL-03: \texttt{gxl.ebnf} \S7.)
\end{quote}

(Source: \texttt{gxl.ebnf} \S5.2, \S5.4.)

% ─────────────────────────────────────────────────────────────────────────────
\subsection{Error Classes}
\label{sec:gxl:errors}
% ─────────────────────────────────────────────────────────────────────────────

GXL defines two timing categories for errors:
\emph{parse errors} are raised at runbook load/plan time, before any step
executes;
\emph{evaluation errors} are raised at runtime when the expression is
evaluated for a specific step context.
A runbook containing a parse error MUST NOT enter execution.
An evaluation error causes the affected step to transition to the
\texttt{failed} state; execution halts unless \texttt{continue\_on\_fail:~true}
is set on that step.

Table~\ref{tab:gxl-errors} enumerates every GXL error code.

\begin{table}[ht]
\centering
\begin{tabular}{p{3.0cm}p{3.0cm}p{7.0cm}}
\hline
\textbf{Code} & \textbf{Raised by} & \textbf{Description} \\
\hline
\multicolumn{3}{l}{\textit{Parse errors (raised at load/plan time)}} \\
\hline
\texttt{GXL-PARSE-001} & Lexer/parser & Unexpected token: does not fit any production. \\
\texttt{GXL-PARSE-002} & Lexer & Unterminated string literal: EOF or newline before closing quote. \\
\texttt{GXL-PARSE-003} & Lexer & Invalid number literal: leading zeros (\texttt{007}), \texttt{NaN}, \texttt{Infinity}, or other malformed numeric token. \\
\texttt{GXL-PARSE-004} & Lexer & Invalid escape sequence: unrecognized character after \texttt{\textbackslash}. \\
\texttt{GXL-PARSE-005} & Parser & Unknown namespace: prefix not in the allowlist (\texttt{str}, \texttt{list}, \texttt{regex}), or \texttt{math} referenced in v1. \\
\texttt{GXL-PARSE-006} & Parser & Unknown namespace method: e.g., \texttt{str.explode()} is not a known method of the \texttt{str} namespace. \\
\texttt{GXL-PARSE-007} & Parser/Lexer & Forbidden syntax: \texttt{\&\&}, \texttt{||}, \texttt{!}, \texttt{|}, ternary, infix \texttt{contains}/\texttt{startsWith}/\texttt{matches}, assignment, string concatenation, array/object construction, negative array index, unary \texttt{+}, or bare function call. Does \emph{not} cover keyword-as-identifier (see \texttt{GXL-PARSE-010}). \\
\texttt{GXL-PARSE-008} & Parser & Unmatched parenthesis. \\
\texttt{GXL-PARSE-009} & Parser & Chained comparison: e.g., \texttt{a < b < c}. \\
\texttt{GXL-PARSE-010} & Parser & Keyword used as identifier (e.g., variable named \texttt{and}, \texttt{or}, \texttt{str}). Sole authoritative code for keyword misuse as identifier; \texttt{GXL-PARSE-007} does not apply to this case. \\
\hline
\multicolumn{3}{l}{\textit{Type errors (parse time where inferrable; otherwise evaluation time)}} \\
\hline
\texttt{GXL-TYPE-001} & Type checker & Type mismatch in comparison: comparing values of incompatible types (e.g., string~\texttt{==}~number). \\
\texttt{GXL-TYPE-002} & Type checker & Non-\texttt{bool} in boolean position: a non-bool value used as the final expression result where \texttt{bool} is required. \\
\texttt{GXL-TYPE-003} & Evaluator & Argument type mismatch: wrong type passed to a stdlib function, or arithmetic applied to a non-\texttt{number}. \\
\texttt{GXL-TYPE-004} & Evaluator & Wrong argument count: function called with incorrect number of arguments. \\
\hline
\multicolumn{3}{l}{\textit{Path errors (raised at evaluation time)}} \\
\hline
\texttt{GXL-PATH-001} & Evaluator & Variable or path segment not found: first \texttt{IDENT} not in context map, or intermediate segment missing in object. \\
\texttt{GXL-PATH-002} & Evaluator & Array index out of bounds: index $\geq$ array length. \\
\texttt{GXL-PATH-003} & Evaluator & Index applied to a non-array value: e.g., \texttt{"hello"[0]} or traversal through a null or scalar intermediate. \\
\hline
\multicolumn{3}{l}{\textit{Evaluation errors (raised at evaluation time)}} \\
\hline
\texttt{GXL-EVAL-001} & Evaluator & Variable resolution failure not otherwise covered by \texttt{GXL-PATH-*} codes. \\
\texttt{GXL-EVAL-002} & Evaluator & Division or modulo by zero (\texttt{x/0}, \texttt{x\%0}). \\
\texttt{GXL-EVAL-003} & Evaluator & Invalid regex pattern: \texttt{regex.match} called with a pattern that is not valid RE2 syntax. \\
\texttt{GXL-EVAL-004} & Evaluator & Null in ordered comparison: \texttt{null} used with \texttt{<}, \texttt{>}, \texttt{<=}, or \texttt{>=}. \\
\hline
\end{tabular}
\caption{GXL error class catalog.
  Source: \texttt{gxl.ebnf} \S6.}
\label{tab:gxl-errors}
\end{table}

% ─────────────────────────────────────────────────────────────────────────────
\section{Migration Note}
\label{sec:gxl:migration}
% ─────────────────────────────────────────────────────────────────────────────

\begin{quote}
\noindent\textbf{Informative.}
This section is not normative.
\end{quote}

GXL replaces the prior \texttt{expr-lang/expr} evaluation surface that was
used for all boolean-typed fields (\texttt{when}, \texttt{condition},
\texttt{until}, and related fields) in earlier versions of the GERT schema.
The key authoring differences are:

\begin{itemize}
  \item Logical operators change from \texttt{\&\&}/\texttt{||}/\texttt{!}
    to the keywords \texttt{and}/\texttt{or}/\texttt{not}.
  \item Infix \texttt{contains} (e.g., \texttt{pod\_json contains "Crash"})
    is removed; authors MUST use \texttt{str.contains(pod\_json, "Crash")}.
  \item Member access on variables (e.g., \texttt{config.region}) is now
    formalized as GDP path syntax and is explicitly supported.
  \item Implicit truthiness is removed; a bare variable in \texttt{when:}
    MUST be of type \texttt{bool}.
  \item Silent empty-string substitution for missing variables is removed;
    every missing path is a hard error.
\end{itemize}

A full deprecation appendix, including a migration guide and rewrite rules for
existing fixtures, will be added in a later Stream~B task.
The existing spec sections that still reference \texttt{expr-lang/expr}
syntax (including \texttt{design/gert/sections/03-schema-vnext.tex} and
\texttt{design/gert/sections/02-architecture.tex}) are identified for
rewrite in the Stream~B backlog; they remain in their current form until
those follow-up tasks are completed.
